% Portable theorem guard: define only what the host preamble is missing, so this
% file drops into any of our manuscripts and re-running it is a no-op.
\makeatletter
\@ifundefined{theorem}{\newtheorem{theorem}{Theorem}}{}
\@ifundefined{lemma}{\newtheorem{lemma}[theorem]{Lemma}}{}
\@ifundefined{proposition}{\newtheorem{proposition}[theorem]{Proposition}}{}
\@ifundefined{corollary}{\newtheorem{corollary}[theorem]{Corollary}}{}
\@ifundefined{definition}{\newtheorem{definition}[theorem]{Definition}}{}
\@ifundefined{remark}{\newtheorem{remark}[theorem]{Remark}}{}
\makeatother
% BT1169 layer-45 bridge insert.

\begin{theorem}[The $45$ relation side as three $15$-layers]
\label{thm:relation-fortyfive-three-layers}
The $45$ relation side of Theorem~\ref{thm:boolean-fortyfive-split} admits the
same layer count as the classical cubic-surface $45$-object model.  After a
choice of double-six orientation, that model splits as two oriented-pair
$15$-layers and one pairing $15$-layer:
\[
  45=15+15+15=3\cdot 15.
\]
Thus the Boolean/Clifford relation sector has the correct object-layer form to
be compared with the tritangent $45$: it is three projective $15$-layers over the
same nonzero-mask base.  The theorem proves the layer count and grade-compatible
split; a canonical incidence isomorphism is not asserted here and remains the
next verification target.
\end{theorem}
